\begin{ouverture}[calcul]{3}{Puissances d’exposant positif}{Partie 1 · Nombres et calculs}
{Produits répétés · Puissances de dix · Règles de calcul}
\objectif{Définir et calculer une puissance d’exposant entier positif.}
\objectif{Multiplier des puissances d’un même nombre ou de même exposant.}
\objectif{Résoudre des problèmes de répétition et de grands nombres.}
\prerequis{Carrés et cubes (Cinquième), produits de nombres relatifs.}
\end{ouverture}
\pagestyle{otMain}\essentiel\fontsize{11}{15}\selectfont
\phantomsection\label{cours-3-1}
\essSection{1}{Une écriture des produits répétés}
\begin{definition}Pour un nombre $a$ et un entier $n\ge1$, $a^n$ désigne le produit de $n$ facteurs égaux à $a$ : $a^n=a\times a\times\cdots\times a$. On lit « $a$ exposant $n$ ». Ainsi $3^4=3\times3\times3\times3=81$ et $a^1=a$. Par convention, $a^0=1$ pour $a$ non nul.\end{definition}
\begin{exemple}$(-2)^4=16$, car quatre facteurs négatifs donnent un produit positif, tandis que $-2^4=-16$ : la puissance porte seulement sur 2. De même $(-3)^3=-27$.\end{exemple}
\begin{definition}[Puissances de dix]Pour $n\ge1$, $10^n$ s’écrit 1 suivi de $n$ zéros : $10^6=1\,000\,000$. Multiplier un nombre par $10^n$ décale chacun de ses chiffres de $n$ rangs vers la gauche. Les préfixes des unités correspondent à des puissances de dix :\end{definition}
\begin{center}\renewcommand{\arraystretch}{1.35}\begin{tabular}{c|cccccc}\toprule
Préfixe & nano & micro & milli & kilo & méga & giga\\\midrule
Facteur & $\frac1{10^9}$ & $\frac1{10^6}$ & $\frac1{10^3}$ & $10^3$ & $10^6$ & $10^9$\\\bottomrule
\end{tabular}\end{center}
\begin{attention}L’exposant compte des facteurs : $5^3=125$, alors que $5\times3=15$. Les exposants négatifs et l’écriture scientifique seront étudiés en Troisième.\end{attention}
\clearpage
\phantomsection\label{cours-3-2}
\essSection{2}{Calculer avec des puissances}
\begin{propriete}[Puissances d’un même nombre]Pour des entiers naturels $m$ et $n$ : $a^m\times a^n=a^{m+n}$. Exemple générique : $2^3\times2^4=(2\times2\times2)\times(2\times2\times2\times2)=2^7$ ; on compte $3+4=7$ facteurs.\end{propriete}
\begin{propriete}[Puissances de même exposant]Pour un entier naturel $n$ : $a^n\times b^n=(a\times b)^n$. Exemple : $2^3\times5^3=(2\times5)\times(2\times5)\times(2\times5)=10^3$ ; on regroupe les facteurs par paires.\end{propriete}
\begin{exemple}$10^4\times10^3=10^7$ et $4^5\times25^5=100^5=10^{10}$. Pour comparer $3\times10^5$ et $25\times10^4=2{,}5\times10^5$, on utilise la même puissance de 10 : le premier est le plus grand.\end{exemple}
\begin{attention}Une somme ne se traite pas comme un produit : $2^3+2^4=8+16=24$, alors que $2^3\times2^4=2^7=128$. Le produit $2^3\times3^2$ ne s’écrit pas avec une seule puissance : les bases et les exposants diffèrent.\end{attention}
\begin{savoirfaire}[Revenir à la définition]En cas de doute, écrire tous les facteurs et les compter. Les deux propriétés résument ce comptage ; elles ne s’appliquent qu’à des produits.\end{savoirfaire}

\clearpage\rubrique{Découverte}
\begin{blocLivre}{ex-3-1}\begin{exercice}[1]{Déplier une puissance}
Écrire $5^3$ comme produit et le calculer.
\end{exercice}\end{blocLivre}
\begin{blocLivre}{ex-3-2}\begin{exercice}[1]{Attention aux parenthèses}
Calculer $(-3)^2$, $-3^2$ et $(-3)^3$.
\end{exercice}\end{blocLivre}
\begin{blocLivre}{ex-3-3}\begin{exercice}[1]{Puissances de dix}
Donner l’écriture décimale de $10^5$, $10^1$ et $10^0$.
\end{exercice}\end{blocLivre}
\begin{blocLivre}{ex-3-4}\begin{exercice}[1]{Écrire une puissance}
Écrire sous la forme d’une puissance : $3\times3\times3\times3\times3$ ; $(-2)\times(-2)\times(-2)$ ; $7$.
\end{exercice}\end{blocLivre}
\begin{blocLivre}{ex-3-5}\begin{exercice}[1]{Lire un préfixe}
Convertir 6 km en mètres, 3 Go (gigaoctets) en octets et 4 mg en grammes.
\end{exercice}\end{blocLivre}

\rubrique{Entraînement}
\begin{blocLivre}{ex-3-6}\begin{exercice}[2]{Même nombre}
Écrire sous la forme d’une seule puissance : $2^3\times2^5$ ; $7\times7^4$ ; $10^6\times10^2\times10$.
\end{exercice}\end{blocLivre}
\begin{blocLivre}{ex-3-7}\begin{exercice}[2]{Même exposant}
Écrire sous la forme d’une seule puissance : $2^4\times5^4$ ; $3^3\times(-2)^3$ ; $25^2\times4^2$.
\end{exercice}\end{blocLivre}
\begin{blocLivre}{ex-3-8}\begin{exercice}[2]{Justifier une règle}
En écrivant les facteurs, expliquer pourquoi $4^3\times4^2=4^5$. L’égalité $4^3+4^2=4^5$ est-elle vraie ?
\end{exercice}\end{blocLivre}
\begin{blocLivre}{ex-3-9}\begin{exercice}[2]{Calculer astucieusement}
Calculer sans calculatrice $2^5\times5^5$ puis $8^3\times125^3$.
\end{exercice}\end{blocLivre}
\begin{blocLivre}{ex-3-10}\begin{exercice}[2]{Une erreur à corriger}
Lou écrit $3^2\times3^4=9^6$ et $2^3\times5^2=10^5$. Corriger chaque égalité.
\end{exercice}\end{blocLivre}

\rubrique{Problèmes}
\begin{blocLivre}{ex-3-11}\begin{exercice}[3]{Une quantité qui double}
Une quantité vaut 3 au départ et double à chaque étape. Exprimer sa valeur après six étapes avec une puissance, puis la calculer.
\end{exercice}\end{blocLivre}
\begin{blocLivre}{ex-3-12}\begin{exercice}[3]{Un message transmis}
Une personne envoie un message à 3 personnes : c’est l’étape 1. À chaque étape suivante, chaque nouveau destinataire l’envoie à 3 personnes qui ne l’ont pas encore reçu. Combien de personnes le reçoivent à l’étape 5 ? Combien de personnes connaissent le message après l’étape 5, l’expéditrice comprise ?
\end{exercice}\end{blocLivre}
\begin{blocLivre}{ex-3-13}\begin{exercice}[3]{Des fichiers}
Une archive contient 2 400 images de 2,5 Mo chacune. On utilise les préfixes décimaux. Quelle taille totale en Mo, puis en Go et en octets ?
\end{exercice}\end{blocLivre}
\begin{blocLivre}{ex-3-14}\begin{exercice}[3]{Comparer de grandes distances}
Une distance A vaut $4\times10^8$ m et une distance B vaut $35\times10^7$ m. Laquelle est la plus grande ? Écrire leur somme sous la forme $a\times10^8$ m.
\end{exercice}\end{blocLivre}
\begin{blocLivre}{ex-3-15}\begin{exercice}[3]{Une somme à ne pas confondre}
Calculer $A=3{,}2\times10^5+8\times10^4$, puis $B=3{,}2\times10^5\times8\times10^4$. Expliquer la différence.
\end{exercice}\end{blocLivre}

\clearpage
